% VI27-DETAILED-PROBLEM-11
\begin{problem}[Dominance sends generic to generic]\label{prob:vi27-p11}
Let \(f:X\to Y\) be dominant between integral schemes. Show \(f(\eta_X)=\eta_Y\).
\par\smallskip\noindent\textit{Provenance.} Canonical VI/27 dossier.
\end{problem}

\ifdefined\FullProblemDossiers
\begin{solution}
Since \(\overline{\{\eta_X\}}=X\), continuity gives
\[
f(X)=f(\overline{\{\eta_X\}})\subseteq\overline{\{f(\eta_X)\}}.
\]
Taking closures and using dominance yields
\[
Y=\overline{f(X)}\subseteq\overline{\{f(\eta_X)\}}\subseteq Y.
\]
Thus \(\overline{\{f(\eta_X)\}}=Y\), so \(f(\eta_X)\) is a generic point of \(Y\). The generic point of an
integral scheme is unique, hence \(f(\eta_X)=\eta_Y\).
\end{solution}
\fi
